\section{La naturaleza de la salida estructurada y las llamadas a herramientas}

\subsection{Por qué se necesita la salida estructurada}

Al construir Agent Workflows complejos, los módulos necesitan un \textbf{acuerdo de interfaz} explícito: la salida de una etapa es la entrada de otra. La salida estructurada resuelve tres problemas centrales:

\begin{enumerate}
    \item No se necesita acordar el Output Format a nivel de Prompt
    \item No se necesita analizar manualmente JSON complejo
    \item La información de Function / Tools tampoco necesita colocarse explícitamente en el Prompt
\end{enumerate}

\begin{importantbox}{Salida estructurada = SDD en versión débil}
La salida estructurada implementa, en cierto sentido, una \textbf{versión débil del Specification Driven Development}: al definir el formato de entrada y salida entre los Módulos, simplifica enormemente el trabajo de análisis y adaptación. La razón subyacente es que los modelos de lenguaje grandes actuales ya están muy bien entrenados en la salida estructurada.
\end{importantbox}

\subsection{Resumen de la sección}

La salida estructurada y las llamadas a herramientas son la infraestructura del desarrollo de Agentes: mediante un acuerdo de formato a nivel de API, se omite una gran cantidad de trabajo manual de diseño de Prompt y análisis de la salida.

